\chapter{知识的获取概述}
\label{chap:knowledge-acquisition-overview}

拍下一批花卉照片，并不等于已经拥有了训练识别模型所需的知识。照片记录颜色、纹理和形状，却未必说明物种名称；同一朵花从不同角度拍摄得到的两张照片，也不会自动附带“来自同一对象”的关系。即使每张照片都有类别，模型仍可能不知道花朵位于哪里，或者哪些局部变化不应改变类别判断。问题因而不只是怎样收集更多数据，而是怎样取得能够解释数据、约束学习并支持目标任务的答案、关系、约束和训练目标。

本部分把\term{知识获取}（knowledge acquisition）限定为：从人工、规则、模型或数据本身等可用来源中，取得可记录、可追溯并能用于后续训练的信息。这个操作性定义不把所有信息都称为知识，也不把已经压入模型参数的能力与下一次训练可读取的外部记录混为一谈。本章先用花卉识别任务界定“缺少什么”和“获得了什么”，再说明六类研究内容怎样衔接，最后沿代表性工作梳理研究问题如何从逐例提供答案扩展到多来源、内部目标和模型生成材料。

\section{知识获取问题}
\label{sec:knowledge-acquisition-problem}

先看四条形式不同的记录。“照片甲的物种是鸢尾”给出一个类别答案；“照片乙与照片丙来自同一株花”给出样本关系；“玫瑰与鸢尾在当前单标签体系中互斥”给出类别约束；“根据未遮蔽区域恢复被遮蔽花瓣”则给出一个训练目标。它们分别支持分类损失、关系一致性、受约束预测和表示学习，不能因为最终都存进训练文件就把它们视为同一种标签。

\begin{definition}[知识记录]
  \label{def:knowledge-acquisition-record}
  在本部分中，一条\term{知识记录}（knowledge record）包含六项：
  \[
    r=(o,s,c,\Omega,z_{\mathrm{sem}},z_{\mathrm{ver}}).
  \]
  其中$o$是记录针对的对象，$s$是来源、取得条件及执行状态，$c$是答案、关系、约束或目标等内容，$\Omega$是内容适用的数据范围与任务范围。语义状态$z_{\mathrm{sem}}$区分已给出、未知与不适用，验证状态$z_{\mathrm{ver}}$区分候选、已核验与已拒绝。只有对象、内容和适用范围足以让训练过程正确解释，来源与两类状态足以让后续过程追溯和检查时，这条记录才具有明确的使用语义。
\end{definition}

这里的\term{知识}是能够解释观测、约束学习或支持任务的上位概念；\term{候选知识}是已经形成记录、但尚未确认能否用于既定范围的知识。类别答案写入某个样本时形成\term{标签}，标签只是知识记录的一种具体形式；标签、关系、约束或恢复目标进入训练过程并影响参数更新时，才承担\term{监督信号}的作用。同一条记录可以同时是候选知识、标签和监督信号，但三个名称分别强调验证状态、记录形式和训练用途，不能互相替代。

这个表示还刻意保留了几个容易被文件格式抹平的区别。\term{未知}表示当前证据不足以确定内容，\term{不适用}表示任务范围本来就不包含该对象；\term{弃权}则是来源或融合系统在当前条件下不输出可用判断的动作。未请求与执行失败属于来源执行状态，融合后证据仍不能唯一决定时记为未决。某位植物学家可以因证据不足而弃权，规则也可以因花朵被遮挡而弃权；复核者随后能够确认这次动作符合规范，但不能据此把未知内容改成已知。“不是玫瑰”才是否定内容。若把这些情况都编码为类别0，模型得到的不是更多知识，而是含义混乱的监督信号。

图~\ref{fig:knowledge-acquisition-observation-to-knowledge}把同一批花卉照片放在三个层次中观察。原始像素是观测；由人工、规则、模型或数据变换形成的答案、关系、约束与目标先是候选知识；训练过程再依据其语义选择相应的损失或约束。图中的“待核验”尤其重要：一个目标能够被自动构造，只说明形式上得到了记录，不说明该记录正确，更不说明它有助于最终识别。

\begin{figure*}[htbp]
  \centering
  \begin{tikzpicture}[
  x=1mm,y=1mm,
  every node/.style={font=\figurefont\footnotesize,text=BookInk},
  layer label/.style={anchor=west,font=\figurefont\bfseries\footnotesize,text=BookTeal},
  observation/.style={rectangle,draw=FigureInput,line width=0.9pt,fill=FigureInput!9,minimum width=52mm,minimum height=12mm,align=center},
  candidate/.style={rectangle,rounded corners=1.2mm,draw=BookAmber,dashed,line width=0.8pt,fill=BookAmber!7,text width=30mm,minimum height=17mm,align=center,inner sep=2mm},
  use/.style={rectangle,draw=FigureOutput,line width=0.8pt,fill=FigureOutput!6,text width=30mm,minimum height=16mm,align=center,inner sep=2mm},
  flow/.style={-{Latex[length=2.2mm,width=1.5mm]},draw=BookMuted,line width=0.8pt,rounded corners=1.2mm}
]
  \node[layer label] at (0,7) {原始观测};
  \node[observation] (photos) at (88,3) {同一组花卉照片\\像素、拍摄条件与对象标识};

  \draw[BookRule,line width=0.6pt] (0,-9) -- (164,-9);
  \node[layer label] at (0,-18) {候选知识};
  \node[candidate] (answer) at (28,-27) {类别答案\\“照片甲是鸢尾”\\{\scriptsize 状态：待核验}};
  \node[candidate] (relation) at (68,-27) {对象关系\\“乙、丙来自同一株”\\{\scriptsize 状态：待核验}};
  \node[candidate] (constraint) at (108,-27) {类别约束\\“玫瑰与鸢尾互斥”\\{\scriptsize 状态：待核验}};
  \node[candidate] (restore) at (148,-27) {恢复目标\\“预测被遮蔽花瓣”\\{\scriptsize 效用：待核验}};

  \draw[flow] (photos.south) -- ++(0,-5) -| (answer.north);
  \draw[flow] (photos.south) -- ++(0,-5) -| (relation.north);
  \draw[flow] (photos.south) -- ++(0,-5) -| (constraint.north);
  \draw[flow] (photos.south) -- ++(0,-5) -| (restore.north);

  \draw[BookRule,line width=0.6pt] (0,-40) -- (164,-40);
  \node[layer label] at (0,-49) {训练使用};
  \node[use] (class loss) at (28,-57) {监督分类损失\\学习物种答案};
  \node[use] (relation loss) at (68,-57) {关系一致性目标\\拉近同一对象表示};
  \node[use] (constrained prediction) at (108,-57) {受约束预测\\排除不合法组合};
  \node[use] (reconstruction) at (148,-57) {内容恢复目标\\学习图像表示};

  \draw[flow,draw=BookTeal] (answer.south) -- (class loss.north);
  \draw[flow,draw=BookTeal] (relation.south) -- (relation loss.north);
  \draw[flow,draw=BookTeal] (constraint.south) -- (constrained prediction.north);
  \draw[flow,draw=BookTeal] (restore.south) -- (reconstruction.north);
\end{tikzpicture}

  \caption{从原始观测到可用知识的概念分层。上层花卉照片只是观测；中层四类候选记录分别表达答案、关系、约束和恢复目标，并分开标明语义状态与验证状态；下层说明训练时如何使用这些记录。虚线边框表示自动或人工得到的记录仍需核验。图中内容为教学构造，不表达实际花卉识别效果。}
  \label{fig:knowledge-acquisition-observation-to-knowledge}
\end{figure*}

明确记录形式以后，还要把“知识不足”拆成可以处理的问题。如果目标地区连原始照片都没有，首先缺少的是观测，需要先采集数据；本章讨论的知识缺口，是已有观测还缺少解释依据、关系、约束或训练目标，或者这些记录的覆盖与质量不足。没有任何稀有物种记录属于内容缺失；常见物种照片很多而稀有物种只出现一次，属于覆盖不足；只有整图类别而没有花朵位置，属于粒度不足。已有答案也可能错误，或者因分类体系、拍摄设备与目标地区变化而过时。多拍同一盆常见花只能增加相似观测，不能自动补上稀有物种的判断依据；把旧标签复制到新类别体系，也不能解决时效失配。表~\ref{tab:knowledge-acquisition-gap-types}把这五类缺口分别连接到需要补充的内容和可观察的验收依据。

\begin{table*}[htbp]
  \centering
  \small
  \setlength{\tabcolsep}{3pt}
  \begin{tabularx}{\linewidth}{@{}>{\setlength{\parindent}{0pt}\RaggedRight\arraybackslash}p{25mm}>{\setlength{\parindent}{0pt}\RaggedRight\arraybackslash}p{43mm}>{\setlength{\parindent}{0pt}\RaggedRight\arraybackslash}p{45mm}>{\setlength{\parindent}{0pt}\RaggedRight\arraybackslash}X@{}}
    \toprule
    \textbf{不足类型} & \textbf{花卉识别中的表现} & \textbf{需要补充或改变的内容} & \textbf{可观察的验收依据} \\
    \midrule
    内容缺失 & 目标地区的某个稀有物种没有任何已确认记录 & 取得该物种的类别依据、适用范围和来源信息 & 目标类别出现已核验记录，且对象与范围字段完整 \\
    覆盖不足 & 常见物种有大量近重复照片，稀有物种或特殊光照几乎未覆盖 & 增加困难物种、环境和外观变化的记录，而非继续复制高频区域 & 预先定义的物种与条件分层均达到最低覆盖要求 \\
    粒度不足 & 每张图只有整图类别，无法回答花朵位置或同一对象关系 & 增加区域、实例或成对关系等与目标问题匹配的细粒度内容 & 在留出的区域或关系样本上能够按相同语义复核 \\
    内容错误 & 相似物种被混淆，或自动答案把背景叶片当成花朵 & 复核并修正内容，同时保留原来源、错误类型和修订轨迹 & 独立复核中的一致性提高，已知错误得到闭环处理 \\
    时效失配 & 分类体系已调整，旧记录仍使用废弃名称；部署地区的物种分布也已变化 & 更新类别映射、适用范围或重新获取受影响记录 & 新体系下无悬空类别，目标时间窗口的抽查达到要求 \\
    \bottomrule
  \end{tabularx}
  \caption{知识不足的类型与相应获取需求。同样是“数据不够”，缺少稀有物种记录与缺少花朵位置需要补充的内容并不相同；验收依据因而必须在获取前随目标一起确定。}
  \label{tab:knowledge-acquisition-gap-types}
\end{table*}

取得记录之后，可以从三个层次判断进展。\term{形式完整性}（formal completeness）检查必要字段是否存在、状态是否合法；\term{参考一致性}（reference consistency）比较记录与独立人工判断、可信规则或其他参考证据是否相符；\term{下游效用}（downstream utility）检查这些记录进入既定训练过程后，是否改善目标数据上的结果。三者不能互相替代。格式正确的错误标签仍然错误，与一位参考标注者一致也可能只是共享歧义，而训练损失下降还可能来自模型记住了系统性偏差。

因此，候选知识只是“已经取得、尚待确认能否使用”的记录，可靠知识则需要相对于明确范围和证据达到预设要求。人工判断可能因任务歧义而分裂，规则可能在遮挡条件下失效，模型可能复制训练资料中的错误，数据内部目标也可能奖励对实际分类无用的捷径。生成数量和模型置信度只能作为过程线索；来源依据、取得条件、独立核验和下游检验共同决定记录是否可以进入下一步。

\begin{definition}[知识获取目标]
  \label{def:knowledge-acquisition-objective}
  给定目标任务$T$、可行获取方案集合$\mathcal{A}$和资源上限$\bm{B}$，每个方案$a\in\mathcal{A}$指定来源、对象、获取动作与复核方式，产生候选记录集合$\mathcal{K}_{\mathrm{cand}}(a)$及其可交付子集$\mathcal{K}(a)\subseteq\mathcal{K}_{\mathrm{cand}}(a)$。知识获取要选择方案，使交付记录的任务效用$U_T(\mathcal{K}(a))$尽可能高，同时满足可靠性$R(\mathcal{K}(a))$、目标范围覆盖$G(\mathcal{K}(a))$和分项成本约束：
  \begin{equation}
    \begin{aligned}
      \max_{a\in\mathcal{A}}\quad
      & U_T\bigl(\mathcal{K}(a)\bigr) \\
      \text{s.t.}\quad
      & \bm{C}(a)\preceq\bm{B}, \\
      & R\bigl(\mathcal{K}(a)\bigr)\geq r_0,\qquad
        G\bigl(\mathcal{K}(a)\bigr)\geq g_0,
    \end{aligned}
    \label{eq:knowledge-acquisition-objective}
  \end{equation}
  其中$\bm{C}(a)=(C_{\mathrm{dev}}(a),C_{\mathrm{acq}}(a),C_{\mathrm{review}}(a),C_{\mathrm{maint}}(a))$分别记录准备开发、逐项获取、复核修正和持续维护的成本。成本按完整方案核算，因此生成后被拒绝的候选记录同样计入获取与复核投入。不同单位的人工工时、调用费用与计算资源应分别报告，除非已经给出可审计的换算规则。
\end{definition}

式~\eqref{eq:knowledge-acquisition-objective}同时给出两种评价口径：预算固定时比较各方案能达到的覆盖、可靠性和下游效用，目标效果固定时比较达到要求所需的投入。实际开始批量获取前，通常先做小规模试点，确认目标数据范围、允许的未知与错误、需要优先覆盖的困难区域以及停止条件。当前批次可以在验收后交付，但类别体系和数据分布仍会变化；长期方案还要说明何时触发重新抽查、更新或废弃旧记录。

\section{知识获取的研究内容}
\label{sec:knowledge-acquisition-research-scope}

知识获取方案需要连续回答四个问题：知识从哪里来，怎样取得，来自不同来源的记录怎样共同使用，以及怎样判断并改进质量。本部分按主要获取机制区分人工、规则、模型和数据本身四类来源，再把融合与质量改进作为两类后续处理。划分依据不是文件格式，也不是过程是否调用程序；一条类别答案可以来自人、规则或模型，同一种来源也可以产生类别、关系和约束等不同内容。表~\ref{tab:knowledge-acquisition-chapter-map}按这一分工给出后续章节入口。

\begin{table*}[htbp]
  \centering
  \small
  \setlength{\tabcolsep}{4pt}
  \begin{tabularx}{\linewidth}{@{}>{\setlength{\parindent}{0pt}\RaggedRight\arraybackslash}p{31mm}>{\setlength{\parindent}{0pt}\RaggedRight\arraybackslash}X>{\setlength{\parindent}{0pt}\RaggedRight\arraybackslash}p{45mm}@{}}
    \toprule
    \textbf{研究内容} & \textbf{核心问题与主要内容} & \textbf{后续章节} \\
    \midrule
    从人工获取知识 & 怎样把需求变成可回答的任务，并选择交互方式、样本、标注者与投入 & 基于人工的知识获取 \\
    从规则获取知识 & 怎样通过模式匹配、知识关联和约束推理产生有依据的候选判断 & 基于规则的知识获取 \\
    从模型获取知识 & 怎样准备模型、组织输入、计算候选答案，并根据反馈改进获取过程 & 基于模型的知识获取 \\
    从数据本身获取知识 & 怎样利用观测内容、样本关系和数据分组构造训练目标与约束 & 基于数据本身的知识获取 \\
    融合多来源知识 & 怎样描述来源差异，对齐对象与语义，并形成综合结果或安排联合使用 & 知识的融合 \\
    评估与修正知识 & 怎样建立判断依据、估计质量、定位问题，并验证修正是否有效 & 知识的质量评估与修正 \\
    \bottomrule
  \end{tabularx}
  \caption{本部分研究内容与章节分工。前四行按知识的主要来源划分，后两行处理多来源共同使用和质量闭环；它们不是六种互斥的文件格式或算法。}
  \label{tab:knowledge-acquisition-chapter-map}
\end{table*}

四类\term{知识来源}（knowledge source）各有不同的证据基础。人工来源把观察者的判断写入记录，能够处理需要语境和专业经验的边界案例，但任务设计、人员差异和逐项投入会影响结果。规则来源把可说明的模式、知识关联或约束反复应用到适用对象，依据清楚且执行一致，却难以覆盖未预见的例外。模型来源利用已经学得的映射、相似样本或图结构产生候选答案，容易扩大覆盖，但其错误可能继承自训练资料。数据本身则通过内容恢复、跨视图关系或分组结构构造目标，不要求先给出目标类别，却必须检验构造目标与实际任务是否一致。表~\ref{tab:knowledge-acquisition-source-conditions}固定花卉识别这一任务，比较四类来源的准备条件、成本、常见错误和独立核验入口。

\begin{table*}[htbp]
  \centering
  \footnotesize
  \setlength{\tabcolsep}{2pt}
  \begin{tabularx}{\linewidth}{@{}>{\setlength{\parindent}{0pt}\RaggedRight\arraybackslash}p{15mm}>{\setlength{\parindent}{0pt}\RaggedRight\arraybackslash}p{21mm}>{\setlength{\parindent}{0pt}\RaggedRight\arraybackslash}p{22mm}>{\setlength{\parindent}{0pt}\RaggedRight\arraybackslash}p{21mm}>{\setlength{\parindent}{0pt}\RaggedRight\arraybackslash}p{19mm}>{\setlength{\parindent}{0pt}\RaggedRight\arraybackslash}p{27mm}>{\setlength{\parindent}{0pt}\RaggedRight\arraybackslash}X@{}}
    \toprule
    \textbf{来源} & \textbf{准备资源} & \textbf{典型输出} & \textbf{覆盖限制} & \textbf{主要成本} & \textbf{常见错误} & \textbf{独立核验入口} \\
    \midrule
    人工 & 任务说明、界面、人员与示例 & 物种答案、花朵区域、弃权或解释 & 受人员能力、可观察信息与时间限制 & 设计、培训、逐项判断与复核 & 歧义理解、疲劳、系统性偏好、随意作答 & 专家复核、重复标注、金标准样本与留出抽查 \\
    规则 & 词表、物种知识、模式与适用条件 & 匹配结果、知识关联、互斥或层级约束 & 只覆盖已编码条件，例外与概念变化需维护 & 规则开发、冲突处理和版本维护 & 过度匹配、漏掉变体、规则冲突、条件过期 & 反例集、覆盖审计、规则溯源与人工抽查 \\
    模型 & 基础模型、提示或输入组织、推理资源 & 类别候选、区域、相似关系或生成材料 & 依赖模型能力、输入表达与训练分布 & 模型准备、调用、筛选和必要更新 & 高置信错误、偏差复制、提示敏感、来源同质 & 独立标注、异源模型、可验证规则与任务留出集 \\
    数据本身 & 未标注观测、变换方法、配对或分组依据 & 恢复目标、跨视图关系、簇分配或顺序约束 & 只能利用观测中存在且构造过程保留的结构 & 数据处理、增广、配对和训练计算 & 捷径学习、错误配对、变换改变语义、分组混杂 & 目标任务留出集、消融、人工检查与跨分布评价 \\
    \bottomrule
  \end{tabularx}
  \caption{四类知识来源的使用条件。各行都以花卉任务为背景比较准备、输出、覆盖、成本与核验入口；颜色、存储格式或是否自动执行都不是来源分类依据。}
  \label{tab:knowledge-acquisition-source-conditions}
\end{table*}

这些来源可以在一个方案中分工。以稀有花卉识别为例，植物学家先给少量代表照片作答并写出判定规范；规则把物种名录与互斥关系应用到更多记录；已有视觉模型提出候选类别和花朵区域；同一对象的多视图照片还可以构造表示学习目标。可比较的类别判断经过对象与语义对齐后再融合，而恢复目标和关系目标可以作为不同训练要求分别使用。融合不是把所有内容强行投票成一个类别，也不是默认来源越多越可靠：共享同一个基础模型、知识库或采集偏差的来源，可能反复提供同一个错误。

一次性方案通常从试点开始，再进行批量获取、必要的对齐融合、抽查修正和交付验收。持续或迭代方案还要把质量判断送回下一轮，调整需求、来源、样本或预算。每一轮都应记录改变了什么、哪些旧结果仍可保留、哪些记录必须重做，以及满足什么条件才停止。重复调用同一个模型生成同类答案，只增加记录条数而没有增加独立信息时，轮次增长不等于知识增长。

图~\ref{fig:knowledge-acquisition-workflow}用实线表示知识记录如何流动，用虚线表示检查结果怎样调整方案。四类来源并列接入候选集合；只有多份结果需要形成共同判断时才经过对齐与融合，单一来源记录可以直接进入质量判断。质量不达要求时，反馈可能改变来源选择，也可能要求重新获取特定对象，而不是无条件重跑全部步骤。

\begin{figure*}[htbp]
  \centering
  \begin{tikzpicture}[
  x=1mm,y=1mm,
  every node/.style={font=\figurefont\footnotesize,text=BookInk},
  stage/.style={rectangle,draw=BookRule,line width=0.8pt,fill=BookPaper,text width=24mm,minimum height=13mm,align=center,inner sep=2mm},
  source/.style={rectangle,rounded corners=1.2mm,draw=FigureInput,line width=0.8pt,fill=FigureInput!8,text width=14mm,minimum height=10mm,align=center,inner sep=1.5mm},
  fusion/.style={stage,draw=FigureModel,fill=FigureModel!8,text width=20mm},
  check/.style={stage,draw=BookAmber,fill=BookAmber!8,text width=21mm},
  result/.style={stage,draw=FigureOutput,fill=FigureOutput!6,text width=15mm},
  flow/.style={knowledge arrow,draw=BookTeal,line width=0.9pt,rounded corners=1.2mm},
  control/.style={knowledge arrow,draw=BookMuted,line width=0.8pt,rounded corners=1.2mm},
  bus/.style={draw=BookMuted,line width=0.8pt},
  knowledge bus/.style={draw=BookTeal,line width=0.9pt},
  feedback bus/.style={draw=BookAmber,dashed,line width=0.8pt,rounded corners=1.2mm},
  adjust/.style={knowledge arrow,draw=BookAmber,dashed,line width=0.8pt,rounded corners=1.2mm},
  edge label/.style={font=\figurefont\scriptsize,inner sep=1pt,text=BookMuted}
]
  \node[stage] (need) at (15,12) {需求与资源\\范围、效果、预算};
  \node[stage] (select) at (47,12) {来源选择\\对象、机制、投入};

  \node[source] (human) at (83,36) {人工};
  \node[source] (rule) at (83,20) {规则};
  \node[source] (model) at (83,4) {模型};
  \node[source] (data) at (83,-12) {数据构造};
  \node[font=\figurefont\scriptsize\bfseries,text=BookTeal] at (83,46) {执行选定的获取机制};

  \node[stage,draw=FigureModel,fill=FigureModel!8] (candidate) at (126,12) {候选知识\\含来源与状态};
  \node[fusion,text width=24mm] (align) at (83,-46) {对象与语义对齐\\需要时再融合};
  \node[check] (quality) at (126,-46) {质量判断\\覆盖与可靠性\\下游效用};
  \node[result] (usable) at (159,-46) {可用结果\\或停止};

  \draw[control] (need.east) -- (select.west);
  \draw[control] (select.east) -- (67,12);
  \draw[bus] (67,-12) -- (67,36);
  \draw[control] (67,36) -- (human.west);
  \draw[control] (67,20) -- (rule.west);
  \draw[control] (67,4) -- (model.west);
  \draw[control] (67,-12) -- (data.west);

  \draw[knowledge bus] (100,-12) -- (100,36);
  \draw[flow] (human.east) -- (100,36);
  \draw[flow] (rule.east) -- (100,20);
  \draw[flow] (model.east) -- (100,4);
  \draw[flow] (data.east) -- (100,-12);
  \draw[flow] (100,12) -- (candidate.west);

  \draw[flow] ([xshift=-8mm]candidate.south) -- (118,-25) -| (align.north);
  \node[edge label,anchor=south] at (99,-24) {多份可比结果};
  \draw[flow] (align.east) -- (quality.west);
  \draw[flow] (candidate.south) -- (quality.north);
  \node[edge label,anchor=west,align=left] at (129,-16) {单一来源\\或独立目标};
  \draw[flow] (quality.east) -- (usable.west);

  \coordinate (feedback split) at (100,-64);
  \draw[feedback bus] (quality.south) |- (feedback split);
  \draw[adjust] (feedback split) -- (100,-72) -| (select.south);
  \node[edge label,text=BookAmber,anchor=south] at (70,-71) {调整来源选择};
  \draw[adjust] (feedback split) -- (64,-64) -- (64,-23) -- (67,-23) -- (67,-12);
  \node[edge label,text=BookAmber,anchor=south] at (82,-63) {调整获取机制};

  \draw[flow] (4,-82) -- (16,-82);
  \node[anchor=west,font=\figurefont\scriptsize] at (18,-82) {知识流动};
  \draw[adjust] (57,-82) -- (69,-82);
  \node[anchor=west,font=\figurefont\scriptsize] at (71,-82) {质量反馈};
\end{tikzpicture}

  \caption{知识获取、融合与质量改进的一种组织方式。实线表示知识流动，虚线表示根据覆盖、可靠性和效用证据调整方案。人工、规则、模型和数据构造通过独立端口汇入候选集合；单一来源可以绕过融合直接检查，多来源结果只有在对象与语义可对齐时才进入融合。质量反馈分别作用于来源选择与具体获取机制。}
  \label{fig:knowledge-acquisition-workflow}
\end{figure*}

比较获取方案时，应固定目标任务、目标数据范围、评价样本和核算周期，再同时观察覆盖、可靠性、下游表现与总投入。准备成本和持续成本要按来源列明，最终效果要保留独立评价依据。若一边反复查看同一评价集来选择样本、调整规则或改写提示，一边又把该评价集报告为最终证据，就会高估方案效果。跨来源的预算协调由第~\ref{chap:knowledge-fusion}章展开，人工来源内部怎样选样、选人与分配重复标注由第~\ref{chap:knowledge-acquisition-human}章展开；日志、存储和调度的系统实现则不属于本部分的内容边界。

\section{知识获取的研究历史}
\label{sec:knowledge-acquisition-history}

本节回顾的是训练知识的获取问题怎样扩展，而不是完整的人工智能史。为给后续脉络建立有据的背景，本章只取两种早期模式：Rosenblatt在1958年的感知机研究中讨论了怎样依据带答案样本调整连接，MYCIN项目的回顾则记录了怎样把领域专家知识编码为规则并检查、维护知识库。前者展示答案随样本进入学习过程，后者展示判断依据怎样成为可重复执行的外部记录；它们是本章采用的背景例证，不是对监督学习或专家系统起源的唯一归因\scite{rosenblatt1958,buchanan1984rulebased}后续自动化改变了人参与的位置，却没有整体替代人工判断和规则：人仍要界定任务、提供困难答案或核验结果，规则仍能表达可审计的约束。

当可询问对象很多时，研究开始追问“哪一条人工判断最值得购买”。Seung、Opper与Sompolinsky在1992年的委员会查询（query by committee）中，用一组与当前证据相容的候选假设之间的分歧选择查询对象；其查询价值依赖假设空间和委员会能否合理表示不确定性。Blum与Mitchell在1998年的协同训练（co-training）中，则研究在两个视图各自足以预测且满足相应独立性假设时，怎样用少量标签和大量未标注样本相互扩展训练信息。前者改变“问哪个样本”，后者改变“怎样利用已有标签和未标注样本”，两者都不能保证自动增加的标签没有偏差\scite{seung1992querycommittee,blum1998cotraining}

标注规模和来源数量增加后，问题从取得单个答案扩展到解释来源差异。Whitehill等人在2009年的模型中联合推断图像标签、标注者能力和样本难度，说明简单多数票会忽略谁在什么样本上更可靠。Ratner等人在2017年发布的Snorkel让专家把启发式知识写成标注函数，再用生成式来源模型估计这些函数的准确性和相关性。两项工作分别从人员与程序化规则进入多来源建模，但共同前提是观测到的分歧确实含有足以区分来源与真值的信息；来源高度相关时，“多人同意”或“多条规则同意”都可能夸大证据\scite{whitehill2009crowdlabels,ratner2017}

\Needspace{18\baselineskip}
另一条路线从观测内部构造目标。BERT的遮蔽语言模型先遮住部分词元，再要求模型根据上下文恢复原内容，使没有下游任务标签的文本也能参与表示学习；该工作于2018年发布预印本，2019年发表于NAACL。这里增加的是预训练目标及由此学得的表示，不是目标任务的正确答案。Bommasani等人在2021年的\term{基础模型}（foundation model）报告中进一步把“大规模宽数据训练、再适配多种下游任务”的模式作为共同研究对象，同时强调下游系统会继承基础模型的缺陷。因而，构造内部目标、学得可迁移表示和获得花卉类别答案是三个需要分别验证的结果\scite{devlin2019,bommasani2021foundation}

生成模型随后直接参与构造训练材料。Self-Instruct让语言模型生成指令、输入与输出，过滤无效或相似项目后再用于微调，展示了“共同生成新任务材料”与“给已有样本预测答案”的区别。DeepSeek-R1报告中的训练流程又组合了模型生成、拒绝采样、可验证任务的规则奖励以及面向一般场景的奖励模型。两项工作都把生成与筛选连成过程，但筛选器仍有覆盖范围，生成材料也可能继承生成模型的盲点；本章只讨论训练知识从哪里来和如何检查，不展开强化学习与模型训练算法\scite{wang2023selfinstruct,deepseek2025r1}

近期研究进一步把选择、生成和核验共同设计。2026年发表于TACL的ActiveLLM使用大语言模型在文本少样本场景中选择待标注实例，回应传统查询策略的冷启动困难，但其证据限定在论文所研究的模型、数据集和选择协议。弱到强泛化研究把弱模型生成的标签用于训练更强模型，观察到强学习者在若干任务上可以超过弱监督者，却仍未恢复强模型的全部能力，不能据此认为弱来源普遍足以监督更强系统\scite{bayer2026activellm,burns2023weaktostrong}

模型输出反复回流还带来新的覆盖风险。Shumailov等人研究了递归使用模型生成数据的情形，在其理论与实验设置中观察到原分布尾部逐代消失的\term{模型坍塌}（model collapse）现象。这个结论针对不加区分地递归使用生成内容的条件，不等于任何合成数据都会导致退化；它提醒获取方案保留真实数据、来源谱系与独立证据，并检查新一轮材料是否真的增加了信息\scite{shumailov2024collapse}

图~\ref{fig:knowledge-acquisition-history}把这些工作放回四条相互重叠的问题线。时间轴只规定先后位置，泳道说明本书用什么问题联系它们；同一工作可能同时影响多条路线，连线也不表示后一个方法由前一个方法直接演化而来。人工标注和专家规则作为持续背景保留，不为缺少统一起点的实践强行指定年份。

\begin{figure*}[htbp]
  \centering
  \begin{tikzpicture}[
  x=1mm,y=1mm,
  every node/.style={font=\figurefont\scriptsize,text=BookInk},
  lane label/.style={anchor=west,font=\figurefont\bfseries\scriptsize,text width=20mm,align=left},
  timeline/.style={draw=BookRule,line width=0.8pt},
  event/.style={rectangle,rounded corners=1mm,line width=0.8pt,minimum height=9mm,align=center,inner sep=1.2mm},
  selection event/.style={event,draw=FigureInput,fill=FigureInput!8,text width=20mm},
  source event/.style={event,draw=FigureModel,fill=FigureModel!8,text width=20mm},
  data event/.style={event,draw=BookTeal,fill=BookTeal!7,text width=20mm},
  generation event/.style={event,draw=BookAmber,fill=BookAmber!8,text width=20mm},
  stem/.style={draw=BookMuted,line width=0.6pt},
  time arrow/.style={knowledge arrow,draw=BookMuted,line width=0.8pt}
]
  \fill[BookPaper] (0,-45) rectangle (22,43);
  \draw[BookRule,line width=0.6pt] (22,-45) -- (22,43);

  \node[lane label,text=FigureInput] at (1,28) {选择人工判断};
  \node[lane label,text=FigureModel] at (1,9) {组织多来源与规则};
  \node[lane label,text=BookTeal] at (1,-10) {利用数据内部结构};
  \node[lane label,text=BookAmber] at (1,-31) {模型生成与核验};

  \node[anchor=south west,font=\figurefont\scriptsize\bfseries,text=BookTeal] at (24,46) {持续背景：人工逐例标注、专家规则、任务设计与核验};
  \draw[time arrow] (24,42) -- (166,42) node[above left,font=\figurefont\scriptsize] {年份：位置保持先后，间距不按比例};
  \draw[timeline] (24,28) -- (166,28);
  \draw[timeline] (24,9) -- (166,9);
  \draw[timeline] (24,-10) -- (166,-10);
  \draw[timeline] (24,-31) -- (166,-31);

  \node[selection event] (qbc) at (32,28) {1992\\委员会查询};
  \node[selection event,text width=17mm] (activellm) at (159,28) {2026\\ActiveLLM};

  \node[source event,text width=22mm] (whitehill) at (64,9) {2009\\Whitehill等};
  \node[source event] (snorkel) at (88,9) {2017\\Snorkel};

  \node[data event] (cotraining) at (49,-10) {1998\\协同训练};
  \node[data event,text width=22mm] (bert) at (95,-10) {2018/2019\\BERT};
  \node[data event,text width=22mm] (foundation) at (120,-10) {2021\\基础模型报告};

  \fill[BookAmber] (131,-31) circle (1.2pt);
  \draw[stem] (131,-31) -- (131,-26);
  \node[generation event,anchor=south] (self instruct) at (131,-26) {2023\\Self-Instruct};
  \draw[stem] (131,-31) -- (131,-36);
  \node[generation event,anchor=north,text width=24mm] (weak strong) at (131,-36) {2023预印本／2024发表\\弱到强泛化};
  \node[generation event,text width=16mm] (collapse) at (142,-31) {2024\\递归生成退化};
  \fill[BookAmber] (153,-31) circle (1.2pt);
  \draw[stem] (153,-31) -- (153,-36);
  \node[generation event,text width=18mm,anchor=north] (deepseek) at (153,-36) {2025\\DeepSeek-R1};

\end{tikzpicture}

  \caption{知识获取研究的问题扩展。横向位置按代表工作的年份排列，四条泳道分别表示选择人工判断、组织多来源与规则知识、利用数据内部结构、模型生成与核验。连线只表示本章解释的问题联系，不表示直接继承、因果关系或后来路线取代早期路线；“2018/2019”分别指BERT预印本与正式发表。}
  \label{fig:knowledge-acquisition-history}
\end{figure*}

表~\ref{tab:knowledge-acquisition-history-works}沿研究问题、获取机制、信息增量和验证边界横向比较图中的代表工作。这个比较不把跨任务性能压成统一排名，也不以发表年份推断方法优劣。

\begin{table*}[htbp]
  \centering
  \footnotesize
  \setlength{\tabcolsep}{2.5pt}
  \begin{tabularx}{\linewidth}{@{}>{\setlength{\parindent}{0pt}\RaggedRight\arraybackslash}p{28mm}>{\setlength{\parindent}{0pt}\RaggedRight\arraybackslash}p{29mm}>{\setlength{\parindent}{0pt}\RaggedRight\arraybackslash}p{34mm}>{\setlength{\parindent}{0pt}\RaggedRight\arraybackslash}p{34mm}>{\setlength{\parindent}{0pt}\RaggedRight\arraybackslash}X@{}}
    \toprule
    \textbf{代表工作} & \textbf{研究问题} & \textbf{来源或获取机制} & \textbf{增加的可用信息} & \textbf{主要假设、验证边界与后续入口} \\
    \midrule
    委员会查询（1992） & 哪些未标注样本值得询问 & 候选假设的预测分歧指导人工查询 & 把有限人工答案集中到当前模型难以确定的样本 & 依赖假设空间与委员会；人工选择见“基于人工的知识获取” \\
    协同训练（1998） & 怎样利用少量标签和大量未标注样本 & 两个视图上的分类器相互提供候选标签 & 扩展可训练的样本标签 & 依赖视图充分性与独立性等条件；数据关系见“基于数据本身的知识获取” \\
    Whitehill等（2009） & 怎样融合能力不同的标注者 & 联合建模标签、标注者能力和样本难度 & 估计综合标签及来源可靠性 & 模型可识别性和设定需检验；人员与融合见“基于人工的知识获取”和“知识的融合” \\
    Snorkel（2017） & 怎样反复执行并组合专家启发式 & 标注函数与生成式来源模型 & 规模化候选标签及来源依赖估计 & 规则覆盖与相关结构可能设错；规则与融合见“基于规则的知识获取”和“知识的融合” \\
    BERT（2018/2019） & 怎样让无下游标签文本参与预训练 & 遮蔽观测并预测原词元 & 内容恢复目标和预训练表示 & 内部目标不等于任务答案；数据构造见“基于数据本身的知识获取” \\
    基础模型报告（2021） & 怎样理解大规模预训练与任务适配 & 宽数据训练后适配多种任务 & 可复用表示与统一问题框架 & 下游会继承上游缺陷；模型来源见“基于模型的知识获取” \\
    Self-Instruct（2023） & 怎样扩大指令训练材料 & 模型生成指令、输入和输出并过滤 & 新的任务描述与训练实例 & 过滤不保证事实正确或分布完整；模型来源见“基于模型的知识获取” \\
    弱到强泛化（2023） & 弱来源能否支持更强学习者 & 弱模型标签监督强模型 & 特定实验中的弱监督训练信号 & 未恢复全部强模型能力；质量边界见“知识的质量评估与修正” \\
    递归生成数据研究（2024） & 模型输出反复回流会发生什么 & 前代模型生成后代训练数据 & 对覆盖退化条件的证据 & 结论依赖递归数据机制；质量与修正见“知识的质量评估与修正” \\
    DeepSeek-R1报告（2025） & 怎样组合生成样本与筛选反馈 & 模型生成、拒绝采样、规则及奖励模型 & 推理与非推理训练材料 & 验证器和生成模型都有覆盖边界；见“基于模型的知识获取”和“知识的质量评估与修正” \\
    ActiveLLM（2026） & 少样本冷启动时怎样选择待标注实例 & 大语言模型根据任务描述选择样本 & 面向人工查询的样本优先级 & 依赖模型和选择协议；人工选样见“基于人工的知识获取” \\
    \bottomrule
  \end{tabularx}
  \caption{代表性研究的问题、获取机制、信息增量与边界。年份采用原始论文的发表年份；BERT同时保留2018年预印本与2019年正式出版口径。表中不比较跨任务性能，也不把代表工作视为相应方向的唯一起点。}
  \label{tab:knowledge-acquisition-history-works}
\end{table*}

\section{本章小结}
\label{sec:knowledge-acquisition-summary}

知识获取不是把照片、标签和损失函数装进同一个文件，而是先说明目标任务缺少什么，再取得含义明确、范围清楚且可追溯的答案、关系、约束或训练目标。花卉识别中的原始照片提供观测，物种答案解释观测，同一对象关系和类别约束补充结构，恢复目标提供另一种训练依据。只有把对象、来源、内容、适用范围、语义状态、执行状态和验证状态保留下来，未知、不适用、弃权、未请求、执行失败与明确否定才不会在训练前被混淆，候选内容也不会被误当成已经核验的结果。

人工、规则、模型和数据本身提供不同证据，融合与质量改进负责让这些证据能够共同使用并经受检查。一个清楚的获取方案应当说明所需知识、可用来源、主要假设、分项成本和验收依据；数量增加、置信度提高或训练损失下降都不能单独证明知识可靠且有用。固定预算下的效果与达到固定效果所需的投入，是评价方案的两种互补口径。

历史上的问题从逐例提供答案和编码专家规则，扩展到选择查询、多来源建模、数据内部目标以及模型生成与核验。自动化扩大了可处理规模，也增加了来源依赖、错误累积和分布覆盖退化的风险。后续各章将分别展开四类来源的获取机制，再讨论多来源融合和质量闭环；这些方法可以组合，但每一项新增记录仍要回到“增加了什么信息、依据是什么、适用于哪里”接受检验。
